\documentclass[10pt,a4paper]{amsart}
\usepackage[margin=19mm]{geometry}
\usepackage{amsmath,amssymb,mathtools}
\usepackage[scr=rsfs,cal=euler]{mathalfa}
\usepackage{xfrac}
\usepackage[unicode,hidelinks,bookmarks=false]{hyperref}
\newcommand{\ie}{i.e.\ }
\newcommand{\cf}{cf.\ }
\newcommand{\resp}{respectively\ }
\newcommand{\Subsection}{\subsection}
\providecommand{\nobreakdash}{\nobreak\hbox{-}}
\setlength{\emergencystretch}{2em}
\multlinegap0pt
\usepackage{bm}
\usepackage{bbm}
\usepackage{dsfont}
\usepackage{xspace}
\usepackage{cancel}
\usepackage{mathtools}
\usepackage{amsthm}
\makeatletter
\@ifundefined{theorem}{\newtheorem{theorem}{Theorem}}{}
\@ifundefined{lemma}{\newtheorem{lemma}[theorem]{Lemma}}{}
\makeatother
\newcommand{\N}{\mathbb{N}}
\title[On Zeckendorf-Niven numbers and arithmetic progressions]{On Zeckendorf-Niven numbers and arithmetic progressions}
\author{Kelly Lao, Steven J. Miller, Nicholas Rosa, Mark Shiliaev, Garrett Tresch, Tony W. H. Wong, Han Zhang}
\date{Source: arXiv:2606.24006v1 (CC BY 4.0)}
\begin{document}
\maketitle

\noindent\emph{Source and licence.} On Zeckendorf-Niven numbers and arithmetic progressions. Version arXiv:2606.24006v1 (CC BY 4.0), distributed under the Creative Commons Attribution 4.0 International licence, as recorded in the arXiv licence metadata for this exact version and shown on the abstract page https://arxiv.org/abs/2606.24006v1. Full rights notice: licenses/arxiv-2606.24006-CC-BY-4.0.txt.\par\smallskip

\noindent\textit{Editorial scope.} Counted unit: the Theorem whose source label is \texttt{thm:infinLN} in the article (source0.tex lines 277--279, source label \texttt{thm:infinLN}), delivered together with the article's own complete original proof of it (source0.tex lines 280--320, one \texttt{proof} environment of 41 lines). No mathematics is written, repaired or re-derived: statement and proof are the article's own text line for line. Required context retained uncounted: thm:infinZN (lines 204--206); lem:piL (lines 254--256); lem:prescribedsL (lines 260--262). Every definition, display and label of those lines is retained unchanged. Omitted from this package and not redistributed: 1--203, 207--253, 257--259, 263--276, 321--516. Nothing inside a retained excerpt is omitted or altered: every retained body is delivered exactly as the article's lines are written, so no omission receipt of any kind accompanies this package. Citations: the retained excerpts carry no citation command at all, so no bibliography entry of the article is transcribed and no reference is redistributed. Reference adaptation: none was needed and none was made; every reference inside the delivered unit resolves inside it, so no unresolved reference or citation is produced. Printed slips kept literal: no printed word, symbol, hypothesis or number of the counted statement or of its proof was altered. Layout and class: the article's own class is \texttt{article} with packages including amsmath, amsthm, amsfonts, amssymb, mathtools, fullpage, multirow, array, bbm, authblk, verbatim, pgf, tikz, pgfplots; the wrapper is the lane's pinned \texttt{amsart} editorial wrapper at 10pt on A4, which restates the theorem-like environments the retained text needs and adds the article's own macro definitions that the retained text needs (\texttt{\textbackslash N}), with extra wrapper packages (bm, bbm, dsfont, xspace, cancel, mathtools). The delivered numbering is the unit's own. Assets: no retained span contains a figure environment, a graphics-inclusion command, a PDF-page inclusion, a table or a TikZ picture; no image file is copied into this package, the package declares no asset dependency and no image file is redistributed with it. Licence: version arXiv:2606.24006v1 of this article is distributed under the Creative Commons Attribution 4.0 International licence, as recorded in the arXiv licence metadata for this exact version and shown on the abstract page https://arxiv.org/abs/2606.24006v1. A full rights notice is delivered at licenses/arxiv-2606.24006-CC-BY-4.0.txt. Characters: every retained excerpt is pure ASCII, so the delivered engine, XeLaTeX with its default Latin Modern text font, reports no missing character.
\begin{theorem}\label{thm:infinZN}
Every arithmetic progression contains infinitely many Zeckendorf-Niven numbers.
\end{theorem}

\begin{lemma}\label{lem:piL}
For every prime $p$ such that $p\equiv-1\pmod{10}$, we have $\pi_L(p)\mid p-1$.
\end{lemma}

\begin{lemma}\label{lem:prescribedsL}
Let $(a+kd)_{k=0}^\infty$ be a given arithmetic progression. Then for every integer $m>s_L(a)$, there exists a term $N$ in the arithmetic progression $(a+kd)_{k=0}^\infty$ such that $s_L(N)=m$.
\end{lemma}

\begin{theorem}\label{thm:infinLN}
Every arithmetic progression contains infinitely many Lucas-Niven numbers.
\end{theorem}

\begin{proof}
With the same logic presented in the proof of Theorem~\ref{thm:infinZN}, it suffices to show that for every $a,d\in\N$, there exists a Lucas-Niven number in the arithmetic progression $(a+kd)_{k=0}^\infty$. By Dirichlet's Theorem, there are infinitely many primes $p$ such that $p\equiv-1\pmod{10\pi_L(d)}$. Let $p>s_L(a)$ be such a prime. Note that $\gcd(p-1,\pi_L(d))\leq2$, which implies that $\gcd(\pi_L(p),\pi_L(d))\leq2$ by Lemma~\ref{lem:piL}. Since $p>s_L(a)$, there exists a term $N$ in the arithmetic progression $(a+kd)_{k=0}^\infty$ such that $s_L(N)=p$ by Lemma~\ref{lem:prescribedsL}. Let
\[N=\sum_{i=1}^pL_{r_i}\]
be the Lucas decomposition of $N$. We are going to complete this proof by considering the following two cases.\\

\noindent\textit{Case $1$}: $\gcd(\pi_L(p),\pi_L(d))=1$.\\

For each $1\leq i\leq p$, there exists an integer $0\leq m_i<\pi_L(p)$ such that
\[m_i\equiv(-r_i+1)\pi_L(d)^{-1}\pmod{\pi_L(p)},\]
i.e., $\pi_L(p)\mid r_i-1+\pi_L(d)m_i$. Let $t_i=r_i+\pi_L(d)(m_i+i\pi_L(p))$ and define
\[N'=\sum_{i=1}^pL_{t_i}.\]
Note that $\pi_L(p)\geq2$, so the summation provided is the Lucas decomposition of $N'$ since $t_1\geq2$ and $t_{i+1}-t_i=r_{i+1}-r_i+\pi_L(d)\pi_L(p)\geq2$ for all $1\leq i\leq p-1$. Consequently, $s_L(N')=p$. Moreover, $t_i\equiv 1\pmod{\pi_L(p)}$ for all $1\leq i\leq p$, so $N'\equiv\sum_{i=1}^pL_1\equiv0\pmod{p}$, implying that $N'$ is Lucas-Niven. Furthermore, since $\pi_L(d)\mid t_i-r_i$ for all $1\leq i\leq p$, we have $d\mid N'-N$. Therefore, $N'$ is a term in the arithmetic progression $(a+kd)_{k=0}^\infty$.\\

\noindent\textit{Case $2$}: $\gcd(\pi_L(p),\pi_L(d))=2$.\\

Let $E=\{1\leq i\leq p:r_i\text{ is even}\}$ and $O=\{1\leq i\leq p:r_i\text{ is odd}\}$. We consider the following subcases.\\

\noindent\textit{Case $2.1$}: $|E|\geq|O|$.\\

Partition $E$ into $E_1$ and $E_2$ such that $|E_1|=|E|-|O|$ and $|E_2|=|O|$. For each $1\leq i\leq p$, there exists an integer $0\leq m_i<\pi_L(p)/2$ such that
\begin{itemize}
\item $m_i\equiv-(r_i/2)(\pi_L(d)/2)^{-1}\pmod{\pi_L(p)/2}$ if $i\in E_1$, i.e., $\pi_L(p)\mid r_i+\pi_L(d)m_i$;
\item $m_i\equiv-((r_i-2)/2)(\pi_L(d)/2)^{-1}\pmod{\pi_L(p)/2}$ if $i\in E_2$, i.e., $\pi_L(p)\mid r_i-2+\pi_L(d)m_i$; and
\item $m_i\equiv-((r_i-1)/2)(\pi_L(d)/2)^{-1}\pmod{\pi_L(p)/2}$ if $i\in O$, i.e., $\pi_L(p)\mid r_i-1+\pi_L(d)m_i$.
\end{itemize}
Similar to Case $1$, let $t_i=r_i+\pi_L(d)(m_i+i\pi_L(p))$ and define $N'=\sum_{i=1}^pL_{t_i}$. With the same reasoning as in Case $1$, we have $s_L(N')=p$, and we also know that $N'$ is a term in the arithmetic progression $(a+kd)_{k=0}^\infty$. Moreover, $t_i\equiv0\pmod{\pi_L(p)}$ if $i\in E_1$, $t_i\equiv2\pmod{\pi_L(p)}$ if $i\in E_2$, and $t_i\equiv1\pmod{\pi_L(p)}$ if $i\in O$, so
\[N'\equiv\sum_{i\in E_1}L_0+\sum_{i\in E_2}L_2+\sum_{i\in O}L_1\equiv2(|E|-|O|)+3|O|+|O|\equiv2p\equiv0\pmod{p},\]
implying that $N'$ is Lucas-Niven.\\

\noindent\textit{Case $2.2$}: $|E|<|O|$.\\

Partition $O$ into $O_1$ and $O_2$ such that $|O_1|=|O|-|E|$ and $|O_2|=|E|$. For each $1\leq i\leq p$, there exists an integer $0\leq m_i<\pi_L(p)/2$ such that
\begin{itemize}
\item $m_i\equiv-((r_i-2)/2)(\pi_L(d)/2)^{-1}\pmod{\pi_L(p)/2}$ if $i\in E$, i.e., $\pi_L(p)\mid r_i-2+\pi_L(d)m_i$;
\item $m_i\equiv-((r_i-1)/2)(\pi_L(d)/2)^{-1}\pmod{\pi_L(p)/2}$ if $i\in O_1$, i.e., $\pi_L(p)\mid r_i-1+\pi_L(d)m_i$; and
\item $m_i\equiv-((r_i+1)/2)(\pi_L(d)/2)^{-1}\pmod{\pi_L(p)/2}$ if $i\in O_2$, i.e., $\pi_L(p)\mid r_i+1+\pi_L(d)m_i$.
\end{itemize}
Once again, letting $t_i=r_i+\pi_L(d)(m_i+i\pi_L(p))$ and defining $N'=\sum_{i=1}^pL_{t_i}$ give us $s_L(N')=p$ and that $N'$ is a term in the arithmetic progression $(a+kd)_{k=0}^\infty$. Finally, $t_i\equiv2\pmod{\pi_L(p)}$ if $i\in E$, $t_i\equiv1\pmod{\pi_L(p)}$ if $i\in O_1$, and $t_i\equiv-1\pmod{\pi_L(p)}$ if $i\in O_2$, so by defining $L_{-1}=-1$, we have
\[N'\equiv\sum_{i\in E}L_2+\sum_{i\in O_1}L_1+\sum_{i\in O_2}L_{-1}\equiv3|E|+(|O|-|E|)+(-1)|E|\equiv p\equiv0\pmod{p},\]
implying that $N'$ is Lucas-Niven.
\end{proof}

\end{document}
